% Generated draft: requires mathematical and source review.
\documentclass[11pt]{article}
\usepackage[margin=25mm]{geometry}
\usepackage{fontspec}
\setmainfont{Latin Modern Roman}
\usepackage{amsmath,amssymb,mathtools,unicode-math}
\setmathfont{Latin Modern Math}
\usepackage{enumitem,needspace,xurl,hyperref}
\setlength{\parindent}{0pt}
\setlength{\parskip}{5pt}
\setlength{\emergencystretch}{3em}
\newcommand{\N}{\mathbb{N}}
\newcommand{\Z}{\mathbb{Z}}
\newcommand{\Q}{\mathbb{Q}}
\newcommand{\R}{\mathbb{R}}
\newcommand{\C}{\mathbb{C}}
\newcommand{\F}{\mathbb{F}}
\newcommand{\A}{\mathbb{A}}
\newcommand{\PP}{\mathbb{P}}
\newcommand{\E}{\mathbb{E}}
\newcommand{\eps}{\varepsilon}
\newcommand{\dd}{\,\mathrm{d}}
\begin{document}
\section*{UnsolvedMath ID: 3100052}
\textbf{Source problem code:} catalogue entry
\subsection{Definitions and mathematical statement}
Let $p$ be prime and write $\Z/p\Z$ in its usual cyclic order. For $S\subseteq\Z/p\Z$ and a multiplier $m\in\Z/p\Z$, put $mS=\{ms:s\in S\}$. The supplied Alon--Peres question asks for the least $K(p)$ such that
\[
 \forall S\subseteq\Z/p\Z\quad\exists m\quad
       \text{every gap in }mS\text{ has length at most }K(p).
\]
The source does not specify whether $m$ must be nonzero, whether $S$ must be nonempty, or how gaps are measured. In particular, a gap can mean a run of unoccupied residues or the distance between successive occupied residues, with either cyclic or linear endpoint conventions. These choices matter for very small sets and cannot be omitted from a precise extremal problem.

To express the target without selecting those conventions, let $\mathcal C_p$ be the permitted family of subsets, $M_p$ the permitted multipliers, and $G_p(T)$ the specified maximum-gap statistic. The corresponding extremal parameter would be
\[
 K_{\mathcal C,M,G}(p)=
 \max_{S\in\mathcal C_p}\min_{m\in M_p}G_p(mS).
\]
This is a formulation schema: the record does not determine all three inputs. The highlighted intended case restricts $\mathcal C_p$ to subsets containing ``about half'' of the $p$ residues. The exact cardinality or density interval in that phrase is not supplied. Denote the resulting, still convention-dependent half-set parameter by $F(p)$, as in the research summary.

That summary reports a literal answer $K(p)=\lceil p/2\rceil$, and for its intended half-set parameter reports
\[
 (\tfrac12-o(1))\log_2p\ \le F(p)\ \le
 (2\sqrt2+o(1))\sqrt p\qquad(p\longrightarrow\infty
 \text{ through primes}).
\]
It also reports $F(p)\ge c\log p\log\log\log p$ for infinitely many primes $p\equiv3\pmod4$, for some positive constant $c$, using quadratic residues and Graham--Ringrose. These are supplied claims, not verified theorems of this draft: in particular the literal answer is not reconciled with the unspecified empty-set, singleton, and gap conventions.

The stated open quantitative question for the intended version is whether there exists $\delta>0$ such that $F(p)\le p^{1/2-\delta}$, with the range and asymptotic-constant convention to be fixed. The summary says a bound $F(p)\le p^{1/4}/\log p$ would improve the cited Burgess record for gaps between quadratic residues. It also asserts
\[
 F(p)\le\varepsilon p\quad\Longleftrightarrow\quad
 k(\varepsilon,p)\le p/2,
\]
but does not define $k(\varepsilon,p)$ or its parameter range, so this equivalence cannot serve as an independent precise formulation here.
\subsection{Short English statement}
How small can one make the largest gap in a set of residues by multiplying the set modulo a prime, uniformly over all permitted sets? The source emphasizes sets of roughly half the residues and their connection to quadratic residues. Its literal universal question and its intended half-set version must be distinguished until the size, multiplier, and gap conventions are fixed.

The original page was labelled open in its October 2020 snapshot. The registry labels the record partially solved, and its research summary calls the intended question open and the square-root upper bound still best known in 2026. These are supplied assessments; the provenance expressly says later resolutions were not checked item by item. The summary attributes the upper bound to Alon--Peres, Proposition 2.1 (1992), the logarithmic lower bound to Weil, the stronger subsequence bound to quadratic residues and Graham--Ringrose, and the exponent question to Green citing Alon--Peres. Its claimed original contributions and second-moment-barrier analysis are not accompanied by proofs in the record.
\subsection{Sources}
\begin{itemize}
\item Local UnsolvedMath registry, permanent ID 3100052, source code catalogue entry, attributed to Alon and Peres. Dataset curation and metadata: Ulam AI and UnsolvedMath Contributors, CC BY 4.0.
\item Cooper, \emph{Combinatorial Problems I Like}, October 2020 snapshot, Combinatorial Number Theory, source-order bullet 52, \url{https://people.math.sc.edu/cooper/combprob.html}; source extraction by HTML, accessed 2026-07-29.
\item Supplied research summary: Alon--Peres, Proposition 2.1 (1992); Green citing Alon--Peres; Graham--Ringrose; Weil; and Burgess. No titles, complete bibliographic records, or separate locators for these attributions are supplied, and they were not independently checked here.
\item Source provenance marks status evidence NEEDS\_REVIEW and rights NEEDS\_REVIEW: public access was verified, but redistribution permission was not established. This source-specific qualification is retained separately from the dataset licence.
\end{itemize}
\end{document}
